\section*{Chapter \ref{ch:g11:antibiotic-resistance} --- Bacteria and Antibiotic Resistance}

\begin{solution}{exo:g11:antibiotic-resistance:1}
\omterm{def:g11:mutations:mutation}{Mutations} arising during the clone's divisions, and \omterm{prop:g11:antibiotic-resistance:variation}{horizontal gene
transfer} --- \omterm{def:g10:universal-dna:information}{DNA}, usually a plasmid, received from another bacterium.
\end{solution}

\begin{solution}{exo:g11:antibiotic-resistance:2}
A molecule that kills bacteria or stops their growth by acting on a
structure bacteria have and human \omterm{def:g10:cells-common-unit:cell}{cells} lack (wall, bacterial
\omterm{def:g10:cells-common-unit:organelle}{ribosome}, bacterial replication \omterm{def:g11:enzymes-and-phenotype:enzyme}{enzymes}). A cold is viral, and
viruses have none of these targets.
\end{solution}

\begin{solution}{exo:g11:antibiotic-resistance:3}
An \omterm{def:g11:enzymes-and-phenotype:enzyme}{enzyme} that destroys the \omterm{def:g11:antibiotic-resistance:antibiotic}{antibiotic}; a mutated target the drug no
longer binds; a pump expelling the drug (or a wall that excludes it).
\end{solution}

\begin{solution}{exo:g11:antibiotic-resistance:4}
Sensitive (20 above 18), but only just. D, with the widest margin
above its threshold, then A.
\end{solution}

\begin{solution}{exo:g11:antibiotic-resistance:5}
No. On replica plates, resistant colonies appear at the same positions
on every \omterm{def:g11:antibiotic-resistance:antibiotic}{antibiotic} plate, so the resistant \omterm{def:g10:cells-common-unit:cell}{cells} existed in the
original colonies, which never met the drug. The \omterm{def:g11:antibiotic-resistance:antibiotic}{antibiotic} reveals
and selects; it does not induce.
\end{solution}

\begin{solution}{exo:g11:antibiotic-resistance:6}
30 divisions: $2^{30} \approx 10^9$ \omterm{def:g10:cells-common-unit:cell}{cells}. About $10^9$ divisions
produced them: some 10 resistant \omterm{def:g10:cells-common-unit:cell}{cells}, on average --- more if the
\omterm{def:g11:mutations:mutation}{mutation} came early.
\end{solution}

\begin{solution}{exo:g11:antibiotic-resistance:7}
Around day 3, when the sensitive \omterm{def:g10:cells-common-unit:cell}{cells} have fallen to about 0.2\% and
the resistant risen to about 1\%. At day 6 the dashed total is back at
100\%, made almost entirely of resistant \omterm{def:g10:cells-common-unit:cell}{cells} (the sensitive
survivors of day 3 regrew too, but from a smaller base).
\end{solution}

\begin{solution}{exo:g11:antibiotic-resistance:8}
A chromosomal \omterm{def:g10:universal-dna:gene}{gene} passes only to the clone's descendants; a plasmid
also passes horizontally, to unrelated \omterm{def:g10:cells-common-unit:cell}{cells} and other species, and
one transfer gives a \omterm{def:g10:cells-common-unit:cell}{cell} resistance at once without waiting for a
\omterm{def:g11:mutations:mutation}{mutation}.
\end{solution}

\begin{solution}{exo:g11:antibiotic-resistance:9}
At day 3 the sensitive \omterm{def:g10:cells-common-unit:cell}{cells} are down to a fraction of a per cent
while the resistant ones are rising; without the drug all survivors
regrow, and the population that returns is far richer in resistant
\omterm{def:g10:cells-common-unit:cell}{cells} than the original. The relapse no longer responds to the same
drug.
\end{solution}

\begin{solution}{exo:g11:antibiotic-resistance:10}
Hospitals use \omterm{def:g11:antibiotic-resistance:antibiotic}{antibiotics} intensively on many patients close
together: selection is constant and the selected strains are passed
from patient to patient by hands and surfaces.
\end{solution}

\begin{solution}{exo:g11:antibiotic-resistance:11}
Streptomycin acts by binding a ribosomal \omterm{def:g11:gene-expression:protein}{protein}; a changed \omterm{def:g11:gene-expression:protein}{protein} no
longer binds it and translation continues. But the change also
slightly impairs the \omterm{def:g10:cells-common-unit:organelle}{ribosome}'s normal work, so \omterm{def:g11:gene-expression:protein}{protein} synthesis, and
growth, are a little slower --- a cost of resistance.
\end{solution}

\begin{solution}{exo:g11:antibiotic-resistance:12}
$10^{-8} \times 10^{-8} = 10^{-16}$ per division: with $10^{12}$
bacteria, no \omterm{def:g10:cells-common-unit:cell}{cell} is likely to carry both. In tuberculosis, treatment
lasts months and the bacteria number billions, so single-drug
resistance is certain to be selected; two drugs at once leave no
survivor to select.
\end{solution}

\begin{solution}{exo:g11:antibiotic-resistance:13}
Low doses select resistant strains in the animals' gut without
killing the population; the workers acquire them by contact; plasmids
carry the \omterm{def:g10:universal-dna:gene}{genes} into human pathogens and into the environment through
manure. Resistance to the drug, and to related human drugs, rises in
medicine.
\end{solution}

\begin{solution}{exo:g11:antibiotic-resistance:14}
Resistance often costs the bacterium (a slower \omterm{def:g10:cells-common-unit:organelle}{ribosome}, an \omterm{def:g11:enzymes-and-phenotype:enzyme}{enzyme} to
make): without the drug, sensitive \omterm{def:g10:cells-common-unit:cell}{cells} grow slightly faster and
outcompete the resistant ones over many generations. The proportion
falls --- but the resistant \omterm{def:g10:universal-dna:gene}{genes} remain in the population at low
frequency, ready to be selected again.
\end{solution}

\begin{solution}{exo:g11:antibiotic-resistance:15}
Bacteria do not adapt individually. In any large population a few
\omterm{def:g10:cells-common-unit:cell}{cells} carry, by random \omterm{def:g11:mutations:mutation}{mutation} or by a plasmid received from
another \omterm{def:g10:cells-common-unit:cell}{cell}, an \omterm{def:g10:universal-dna:gene}{allele} that makes them resistant. The \omterm{def:g11:antibiotic-resistance:antibiotic}{antibiotic}
kills the others; the resistant \omterm{def:g10:cells-common-unit:cell}{cells} multiply and pass the \omterm{def:g10:universal-dna:gene}{allele}
to their descendants. The population becomes resistant because it
has been sorted, not because its members learnt anything.
\end{solution}

\begin{solution}{pb:g11:antibiotic-resistance:1}
\textbf{1.} About $10^{10}$ divisions (one per \omterm{def:g10:cells-common-unit:cell}{cell} produced).

\textbf{2.} $10^{10} \times 10^{-8} = 100$ resistant \omterm{def:g10:cells-common-unit:cell}{cells}, on average.

\textbf{3.} A hundred \omterm{def:g10:cells-common-unit:cell}{cells} in ten billion are invisible on a plate:
the lawn behaves as sensitive because 99.999999\% of it is.

\textbf{4.} Amoxicillin: sensitive (24 above 17). Erythromycin:
resistant (12 below 20). Penicillin: fully resistant, no inhibition.

\textbf{5.} Penicillin is useless against this strain. The choice
among sensitive drugs also weighs side effects and spectrum; the
narrowest effective drug is preferred to spare the patient's other
bacteria.

\textbf{6.} Four periods of 6 hours per day: $10^{10} \times 0.1^4 =
10^6$ after 1 day; $10^{10} \times 0.1^{12} = 0.01$ after 3 days ---
essentially none.

\textbf{7.} $100 \times 2^{48} \approx \num{2.8e16}$: absurd, more than
the wound could hold. Growth is limited by space, nutrients and the
immune system.

\textbf{8.} From 100 \omterm{def:g10:cells-common-unit:cell}{cells} to $10^{10}$ needs $\log_2 10^8 \approx 27$
doublings, about 13 hours; in practice the space frees over the first
day, so by day 1 to 2 the resistant clone has replaced the sensitive
population.

\textbf{9.} Sensitive: falling by a factor of 10 every 6 hours to
nothing by day 3. Resistant: rising from 100 to fill the space within
a day or two, then flat at $10^{10}$. Total: a dip during day 1, then
back to $10^{10}$, now all resistant.

\textbf{10.} If the immune system clears \omterm{def:g10:cells-common-unit:cell}{cells} faster than the
resistant hundred can multiply, the infection ends despite them; if
the clone outgrows the clearance, the infection becomes resistant.
\omterm{def:g11:antibiotic-resistance:antibiotic}{Antibiotics} work with the immune system, not instead of it.

\textbf{11.} Sensitive after 2 days: $10^{10} \times 0.1^8 = 100$.
Resistant: about $10^{10}$ if the clone has filled the space, or at
least many thousands: the survivors are almost all resistant.

\textbf{12.} Growing at the same rate preserves the proportion:
essentially 100\% resistant.

\textbf{13.} No effect: the resistant clone ignores amoxicillin.

\textbf{14.} The first plate held the original population, 99.999999\%
sensitive; the relapse plate holds the selected clone, all resistant,
so the lawn grows up to the disc.

\textbf{15.} The \omterm{def:g10:universal-dna:gene}{gene} can pass to other species carried by other
patients or staff, so the ward's other pathogens can become resistant
without waiting for their own \omterm{def:g11:mutations:mutation}{mutation}.

\textbf{16.} On the ward, 400 courses a year in a small closed
population sort the bacteria continually, and the selected strains
pass between patients; in the town selection is diluted among many
people and species.

\textbf{17.} Antibiograms: choose a drug the strain cannot resist, so
no clone is selected. Full courses: no interrupted course leaves the
resistant survivors to regrow. Hygiene: no transmission of the
selected clone between patients. Isolation: no horizontal or
vertical spread from carriers.

\textbf{18.} Less selection let sensitive strains, which grow slightly
faster, regain ground, and transmission of resistant ones was cut. It
did not return to 5\% because the resistant strains and their
plasmids persist at low frequency and are re-selected by each
remaining course.

\textbf{19.} Any large bacterial population contains \omterm{def:g10:cells-common-unit:cell}{cells} resistant
to the new drug by chance \omterm{def:g11:mutations:mutation}{mutation}; using it selects them, and in a
few years the resistant strain is common, as with every \omterm{def:g11:antibiotic-resistance:antibiotic}{antibiotic}
introduced so far.

\textbf{20.} About a hundred resistant \omterm{def:g10:cells-common-unit:cell}{cells} among ten billion before
any drug; on the second day of treatment they inherited the wound;
stopping the course early turned the curable infection into a
resistant one that spread.
\end{solution}
